\chapter{The Distance Formula}
\noindent
% --- EDITOR VISIBILITY CHECK ---
% If you can read this, the mobile editor is working.
\colorbox{practicegreen}{\parbox{\dimexpr\linewidth-2\fboxsep}{\centering\small\bfseries\color{white} EXPERIMENTAL LAB: FIELD APPLICATIONS}}
\vspace{15pt}

\begin{enumerate}
    \item \textbf{Laser Rangefinder Calibration:} A surveyor places a laser at origin $(0,0)$ and a reflector at $(15, 20)$. The laser reads $24.8$ units. Calculate the theoretical distance and find the percentage error.
    \item \textbf{The Acoustic Localization:} Two microphones at $M_1(-5, 0)$ and $M_2(5, 0)$ detect sound. A third at $M_3(0, 12)$ detects it later. Find the source $S(0, y)$.
    \item \textbf{Shadow Tracking:} A pole's top is at $(0, 5)$ and its shadow is at $(3, 0)$. When the sun moves, the shadow moves to $(0, 0)$. Calculate the total distance the shadow tip traveled.
    \item \textbf{Tension Wire Stability:} A tower at $(0, 12)$ is secured by wires anchored at $(x, 0)$ and $(-x, 0)$. If the total length of both wires is 26, find $x$.
    \item \textbf{GPS Drift Analysis:} Readings at $(10, 10)$ drift to $(10.1, 9.9)$ and $(9.8, 10.2)$. Find the average distance of these drift points from the center.
\end{enumerate}

\chapter{Coordinate Geometry}
\section{Introduction}
Coordinate geometry, also known as analytic geometry, is the study of geometry using a coordinate system. This contrasts with synthetic geometry.

\section{The Cartesian Plane}
The Cartesian plane is defined by two perpendicular number lines: the x-axis, which is horizontal, and the y-axis, which is vertical.

\chapter{New Chapter}
\noindent
% --- EDITOR VISIBILITY CHECK ---
% If you can read this, the mobile editor is working.
\colorbox{practicegreen}{\parbox{\dimexpr\linewidth-2\fboxsep}{\centering\small\bfseries\color{white} EXPERIMENTAL LAB: ANNIHILATOR OPERATOR}}
\vspace{15pt}

\textbf{Definition:} An \emph{annihilator} is a differential operator that maps a function completely to zero when applied to it.

Let
\[
D=\frac{d}{dx}.
\]

\begin{enumerate}
    \item \textbf{Differentiate the Function:} For the exponential function $e^{2x}$,
    \[
    D(e^{2x})=\frac{d}{dx}(e^{2x})=2e^{2x}.
    \]
    Hence, the derivative is still a multiple of the original function, so $D$ alone does not annihilate it.

    \item \textbf{Apply the Annihilator:} Consider the operator $(D-2)$:
    \[
    (D-2)e^{2x}
    =D(e^{2x})-2e^{2x}
    =2e^{2x}-2e^{2x}
    =0.
    \]
    Therefore, the annihilator of $e^{2x}$ is
    \[
    \boxed{D-2.}
    \]

    \item \textbf{Meme Interpretation:}
    \begin{itemize}
        \item \textbf{Using $D=\dfrac{d}{dx}$:}
        Shoots $e^{2x}$, but it clones itself since
        \[
        D(e^{2x})=2e^{2x}.
        \]

        \item \textbf{Using $(D-2)$:}
        Actually annihilates $e^{2x}$ because
        \[
        (D-2)e^{2x}=0.
        \]
    \end{itemize}
\end{enumerate}


